%% kapa: exam=1 round=32 year=2021 period=1 form=A subject=civil_law subject_ko=민법 seq=36 number=36
%% source: Q-Net 감정평가사 기출문제 (cst00309) artlSeq=5208397 file=extracted/1008/5208397/A형/2021년도 제32회 감정평가사 1차 1교시 A형.pdf
%% conversion: kordoc 4.24.51 -> tools/gen1.py (auto); review=proofread-sonnet; flags=none
\begin{question}{36}
유치권에 관한 설명으로 옳지 않은 것은? (다툼이 있으면 판례에 따름)
\begin{choices}
  \item 건물의 임차인이 임대인에게 지급한 임차보증금반환채권은 그 건물에 관하여 생긴 채권이 아니다.
  \item 임대인이 건물시설을 하지 않아 임차인이 건물을 임차목적대로 사용하지 못하였음을 이유로 하는 손해배상청구권은 그 건물에 관하여 생긴 채권이다.
  \item 수급인의 재료와 노력으로 건축되었고 독립한 건물에 해당되는 기성부분에 대하여는 특별한 사정이 없는 한 수급인은 유치권을 가질 수 없다.
  \item 채권자가 채무자를 직접점유자로 하여 간접점유하는 경우에는 유치권이 성립하지 않는다.
  \item 유치권자가 점유침탈로 유치물의 점유를 상실한 경우, 유치권은 원칙적으로 소멸한다.
\end{choices}
\end{question}
\answer{2}
